\documentclass[12pt]{article}
\usepackage{latexsym, amscd, %graphicx, 
amssymb}
\newcommand{\SECT}[1]
		{\vspace{.3cm}
		$\!\!\!\!\!\!\!\!\!\!\!${\bf {#1}}
		\vspace{.2cm}}

\begin{document}


\begin{center}

{\large \bf Math 152, Fall, 2004, Workshop 3\par}
\vspace{.5cm}
{{\bf Honors Section} 
\par}
\end{center}


\begin{enumerate}


\item Calculate four of the following five indefinite integrals:

\begin{enumerate}

\item ${\displaystyle \int x\cos x^{2} dx}$.

\item  ${\displaystyle \int x^{2}\cos x dx}$.

\item ${\displaystyle \int x^{2}\cos x^{2} dx}$.

\item ${\displaystyle \int x^{2}\cos x^{3} dx}$.

\item ${\displaystyle \int x^{3}\cos x^{2} dx}$.

\end{enumerate}


\item A particle of mass 2 kilogram moves along a line; 
its position relative to the origin at time $t$ second 
is $s(t) = t^{3}$ meters. 

\begin{enumerate}

\item Use Newton's law ($F = ma$) to find the force 
$F = f(t)$ which acts on the body at time $t$, and 
then  find the average value of this force between 
$t = 0$ sec and $t = 2$ sec. This is the time average 
$F_{\mbox{\rm \scriptsize time average}}$ of the force. Draw a graph 
of the function $f(t)$ for $0\le t\le 2$. Find the time $t^{*}$ for 
which $f(t^{*}) = F_{\mbox{\rm \scriptsize time average}}$ 
and give a graphical interpretation. 

\item Find a formula $F = g(s)$ for the force which acts on 
the body when it is at position $s$, and 
then find the averge value of $g(s)$ between $s = s(0)$ and 
$s = s(2)$. This is the distance average 
$F_{\mbox{\rm \scriptsize distance average}}$ of the force 
as the body moves from position $s(0)$ to position $s(2)$. Draw a graph 
of the function $g(s)$ for $s(0)\le  s\le s(2)$. Find the 
position $s^{*}$ for which $g(s^{*}) = F_{\mbox{\rm 
\scriptsize distance average}}$ 
and
give a graphical interpretation. 
(Notice that the time and distance averages of the force are different! )

\end{enumerate}

\item \begin{enumerate}

\item Obtain the reduction formula
$$\int (\ln x)^{n} dx = x(\ln x)^{n}-n \int (\ln x)^{n-1} dx\hspace{6em} (*)$$
using integration by parts with $u=(\ln x)^{n}$ and $dv= dx$.

\item Use the substitution $x=e^{y}$ to replace $(*)$ with a 
formula involving functions of $y$.

\item Give a direct proof of the formula obtained in (b) 
using integration by parts.

\end{enumerate}


\item In Section 2.1, average velocity
was defined as 
$$\mbox{\rm average velocity}=\frac{\mbox{\rm distance traveled}}
{\mbox{\rm time elapsed}}$$
(page 90). This was used to motivate a definition of 
{\it instantaneous velocity} as the derivative of distance 
with respect to time. Thereafter, the word ``velocity'' always
referred to instantaneous velocity. We now have another meaning of the 
word ``average,'' introduced in Section 6.5. Show that the old
meaning of average velocity is the average of a function 
giving velocity as a function of time.

If the velocity is always of the same sign, the function giving distance as 
a function of time has an inverse, so all properties of the motion
can be expressed as a function of distance. It is thus possible 
to consider averaging velocity with respect to distance. These two averages
can be different. Show that this happenes when the motion is given by 
{\it free fall}, with displacement given by $s=\frac{1}{2}gt^{2}$. 

\end{enumerate}


\end{document}



